\section{Additional Experimental Results}\label{sec:extended_results}
The results below are additional to the full published study retained
in Sections~\ref{sec:results}--\ref{sec:comparison}. New campaign outcomes
will be associated with their own inputs, versions, and observation
protocol. They do not silently replace the historical measurements.

\subsection{Code--Data Identification and Functional Outcomes}
\draftnote{Insert the completed comparison, independent reference
annotations if available, candidate/accepted interval results, analysis
time, and functional outcomes after extraction. Keep selection and
architecture explicit.}
\begin{table}[t]
\caption{Additional code--data assessment; entries await finalized
experiments. Use agreement/case-study reporting if independent labels
are unavailable.}
\label{tab:new_data_results}
\centering\small
\begin{tabularx}{\linewidth}{@{}Yllll@{}}
\toprule
Method & Inputs & Precision$_D$ & Recall$_D$ & $F_{1,D}$ \\
\midrule
Symbolic/heuristic & \pending & \pending & \pending & \pending \\
LLM classification & \pending & \pending & \pending & \pending \\
\bottomrule
\end{tabularx}
\end{table}
The final analysis will separate a changed classification from a
changed acceptance decision or functional oracle. Any current
working-rate comparison with the original study will state changes
in corpus selection, preparation, and testing conditions.

\subsection{Assembly Wrapping}
\draftnote{Insert build counts, working rates, size, diversity, and
completed static scans. x86 remains development work until its version
and architecture-specific campaign are fixed.}
\begin{table}[t]
\caption{Additional wrapping outcomes. Report both planned-case and
constructed-output working rates; static results refer to completed scans.}
\label{tab:new_wrap_results}
\centering\small
\begin{tabularx}{\linewidth}{@{}Yllll@{}}
\toprule
Condition & $B/N$ & $W/N$ & $W/B$ & $\overline d$ \\
\midrule
x64 wrapping & \pending & \pending & \pending & \pending \\
x64 wrapping, combined & \pending & \pending & \pending & \pending \\
x86 wrapping & \multicolumn{4}{c}{\dev} \\
x86 wrapping, combined & \multicolumn{4}{c}{\dev} \\
\bottomrule
\end{tabularx}
\end{table}

\subsection{LLM Transformation: Zero-Shot versus Two-Shot}
\draftnote{Insert the final input/run counts, initial and final acceptance,
functional outcomes, static verdicts, normalized sequence similarity,
output uniqueness, attempts, time, and estimated cost. Record the scope
of the pilot and the role of the generation-time verifier.}
\begin{table}[t]
\caption{Additional LLM transformation results. Similarity denotes
the normalized sequence score, distinct from the historical cosine
measure. Pending scans have no measured detection fraction.}
\label{tab:new_llm_results}
\centering\small
\begin{tabularx}{\linewidth}{@{}Yllll@{}}
\toprule
Family / shots & Runs & $W/N$ & $\overline d$ & Similarity \\
\midrule
EQ / 0 & \pending & \pending & \pending & \pending \\
EQ / 2 & \pending & \pending & \pending & \pending \\
NOP / 0 & \pending & \pending & \pending & \pending \\
NOP / 2 & \pending & \pending & \pending & \pending \\
BCF / 0 & \pending & \pending & \pending & \pending \\
BCF / 2 & \pending & \pending & \pending & \pending \\
IBP / 0 & \pending & \pending & \pending & \pending \\
IBP / 2 & \pending & \pending & \pending & \pending \\
\bottomrule
\end{tabularx}
\end{table}
The comparison will interpret representation changes alongside
construction and functional outcomes. Similarity alone does not
establish correctness, and accepted outputs do not characterize the
success rate of first attempts or exhausted runs.
